\section{Whole original function fields from actual projective coordinate ratios}
\label{sec:projective-chart-function-field-generation}

\begin{theorem}[Native projective-chart coefficient algebra]
\label{thm:projective-chart-coefficient-algebra}
\lean{EnvelopeProperty.ProjectiveChartFunctionFieldGeneration.chart_structure, EnvelopeProperty.ProjectiveChartFunctionFieldGeneration.sectionsEquiv, EnvelopeProperty.ProjectiveChartFunctionFieldGeneration.sectionsEquiv_coordinate}
\leanok
\uses{thm:scheme-affine-coefficient-equivalences}
\source{bj22-trivval:sec:tc}
The actual polynomial coordinate chart of the original
$\mathbf P^m_k$ has its ordinary polynomial coordinate algebra
as a $k$-algebra for the original projective coefficient map.
The unchanged coordinate-ring isomorphism sends every original
normalized coordinate to the actual section $X_j/X_i$ on the
whole original chart.
\end{theorem}

\begin{proof}
The actual affine coordinate-chart scheme isomorphism has its
original coefficient structure square. Cancel its inverse to
identify the coefficient projection of the original open chart
with the projection of its polynomial spectrum. The preceding
native coefficient-isomorphism theorem therefore makes its
actual section isomorphism a $k$-algebra isomorphism. Compose
with the actual original open top-section isomorphism.
The ring isomorphism is the original dehomogenization and
homogeneous-localization isomorphism followed by the actual
Proj section map. Its image on a normalized coordinate is the
literal degree-one fraction $X_j/X_i$, as asserted.
\end{proof}

\begin{theorem}[Every original closed embedding internally generates its whole original function field]
\label{thm:projective-chart-function-field-generation}
\lean{EnvelopeProperty.ProjectiveChartFunctionFieldGeneration.coordinate, EnvelopeProperty.ProjectiveChartFunctionFieldGeneration.coordinate_self, EnvelopeProperty.ProjectiveChartFunctionFieldGeneration.adjoin_coordinates_eq_top}
\leanok
\uses{thm:projective-chart-coefficient-algebra,thm:affine-function-field-generation,thm:scheme-function-field-pullback-coefficients}
\source{bj22-trivval:sec:tc}
Let $f:X\hookrightarrow\mathbf P^m_k$ be an actual closed
embedding of an integral original $X$. Choose any original
coordinate chart with nonempty inverse image. The actual
pullbacks of $X_j/X_i$, followed by the original generic
germ into the native $K(X)$, have pivot one and generate all
of that same original $K(X)$ over the original coefficient
field. This field-generation conclusion is derived from the
closed embedding and actual coordinate geometry.
\end{theorem}

\begin{proof}
The pivot section $X_i/X_i$ is one, and actual pullback and
the generic germ preserve one. Form the intermediate field
generated by all the original pulled coordinate ratios.
The native chart polynomial algebra is generated by the
nonpivot ratios. The original pullback followed by the
original generic germ is a $k$-algebra homomorphism, so
polynomial induction shows that this field contains the
germ of every actual original chart section: it contains
all coordinate images, all original scalar images, their
sums and their products. Actual closed immersion makes the
whole affine-chart section map surjective. The preceding
affine generation theorem therefore says these same chart
section germs generate all of the native $K(X)$, forcing
the coordinate-generated intermediate field to be all of it.
\end{proof}

\begin{theorem}[Actual homogeneous model over the entire original function field]
\label{thm:projective-chart-whole-native-field-model}
\lean{EnvelopeProperty.ProjectiveChartFunctionFieldGeneration.fullFieldModel}
\leanok
\uses{thm:projective-chart-function-field-generation,def:gl-model}
\source{bj22-trivval:sec:tc}
The same original closed embedding and nonempty coordinate
chart internally construct an integral homogeneous model
of rank one over the entire same native $K(X)$, with actual
coordinates the unchanged pulled ratios $X_j/X_i$ and
actual homogeneous algebra $k[(X_j/X_i)T]$.
\end{theorem}

\begin{proof}
Take the algebra generated by these finitely many original
degree-one monomials. It is finitely generated and graded
because its generators are monomials. Every original
coordinate lies in it, the pivot is nonzero with coefficient
one, and integrality over its own coordinate algebra is
immediate. For the model field-generation condition, divide
each original coordinate monomial by the pivot monomial.
The result is precisely its original coefficient. These
coefficients generate the entire native $K(X)$ by the
preceding theorem, so same-degree quotients generate that
entire field. No field-generation or homogeneous-model
comparison premise is supplied. This does not yet assert
a scheme, analytic-space or numerical-psh equivalence.
\end{proof}

\begin{theorem}[Original native projectivity internally supplies the embedding and generating chart]
\label{thm:native-projective-whole-field-coordinates}
\lean{EnvelopeProperty.ProjectiveChartFunctionFieldGeneration.exists_native_field_coordinates}
\leanok
\uses{thm:projective-chart-function-field-generation}
\source{bj22-trivval:sec:tc}
Every actual integral native-projective $X$ over its original
field internally admits a closed polynomial projective
embedding over its exact original structure map, and an
original nonempty coordinate chart whose actual original
coordinate ratios generate the entire same native $K(X)$.
The original coefficient action of $X$ is retained exactly.
\end{theorem}

\begin{proof}
Actual native projectivity supplies a closed polynomial
projective embedding with its exact structure triangle.
The actual coordinate charts cover projective space, so
one contains the image of the original generic point of
$X$. Its inverse image is nonempty. Apply the preceding
whole-field coordinate theorem to that unchanged embedding
and this actual chart. The exact original structure
triangle identifies its original coefficient action with
the original structure-map action of $X$, giving the stated
whole native field-generation result and pivot one.
\end{proof}
